%!TEX root = ../main.tex
% ANEXO II — Modelo financiero
% Archivo: anexos/Presupuesto_y_Proyecciones_-_Proyecto_Tinku.xlsx

\chapter{Modelo financiero}\label{anx:modelo-financiero}

Los cálculos de la evaluación económica del
capítulo~\ref{cap:economica} se realizan en una planilla de cálculo con
parámetros editables, que se entrega junto con este documento en el
siguiente archivo:

\begin{center}
    \path{Presupuesto_y_Proyecciones_-_Proyecto_Tinku.xlsx}
\end{center}

Todos los montos están expresados en dólares estadounidenses.

\section{Estructura de la planilla}

\begin{table}[H]
    \centering
    \caption{Hojas del modelo financiero}
    \label{tab:hojas-modelo}
    \small
    \begin{tabularx}{\textwidth}{l L}
        \toprule
        \tb{Hoja} & \tb{Contenido} \\
        \midrule
        Inversión Inicial        & Inversión incremental del lanzamiento: horas y tarifa por perfil, dominio, asesoría legal y contingencia, más el costo hundido del \gls{mvp} a título informativo (tabla~\ref{tab:inversion}) \\
        Gastos Post-lanzamiento  & Costos fijos mensuales por año: infraestructura, servidor de aplicación, video, asistente \gls{llm}, honorarios contables, dominio y marketing (tabla~\ref{tab:costos-fijos}) \\
        Valor de Venta           & Comisión, precio por hora, duración promedio de las sesiones, ticket, costo de la pasarela, costo, precio y adopción del resumen opcional, y margen neto por sesión (tabla~\ref{tab:unit-economics}) \\
        Estimación de Demanda    & Alumnos activos por año y escenario, y sesiones anuales resultantes (tabla~\ref{tab:demanda}) \\
        VAN, TIR y Repago        & Tasa de descuento construida (tabla~\ref{tab:tasa}), flujo de fondos de los tres escenarios, \gls{van}, \gls{tir}, repago simple y descontado, y resumen comparativo (tablas~\ref{tab:flujo-caja} y~\ref{tab:indicadores-financieros}) \\
        Gráficos                 & Flujo neto y flujo acumulado por escenario, y \gls{van} comparado \\
        \bottomrule
    \end{tabularx}
\end{table}

\section{Supuestos editables}

Los valores de entrada están separados de los cálculos: al modificar
cualquiera de ellos se recalculan los tres escenarios. Los
principales son la comisión de plataforma, el precio por hora, la
proporción de sesiones de 30 minutos, el precio y la adopción del resumen opcional, las tarifas horarias por perfil,
los alumnos activos por año, las clases por alumno por mes y la tasa de
descuento. Cada supuesto documenta su origen en la columna de
observaciones o en un comentario de la celda.

Los indicadores se calculan con las funciones nativas \texttt{NPV} e
\texttt{IRR}, y los períodos de repago con la interpolación de la
ecuación~\eqref{eq:repago}.
